\documentclass[UTF8,a4paper,11pt,fontset=none]{ctexart}
\usepackage[left=24mm,right=24mm,top=24mm,bottom=24mm,headheight=16pt]{geometry}
\usepackage{fontspec}
% 与三阶段报告相同的开源字体，随 TeX Live 提供。
\setmainfont{texgyretermes-regular.otf}[BoldFont=texgyretermes-bold.otf,ItalicFont=texgyretermes-italic.otf,BoldItalicFont=texgyretermes-bolditalic.otf]
\setsansfont{texgyreheros-regular.otf}[BoldFont=texgyreheros-bold.otf]
\setmonofont{texgyrecursor-regular.otf}[Scale=MatchLowercase]
\setCJKmainfont{FandolSong-Regular.otf}[BoldFont=FandolSong-Bold.otf,ItalicFont=FandolKai-Regular.otf]
\setCJKsansfont{FandolHei-Regular.otf}[BoldFont=FandolHei-Bold.otf]
\setCJKmonofont{FandolFang-Regular.otf}
\usepackage{amsmath,booktabs,longtable,array,seqsplit}
\usepackage{xcolor,graphicx,caption,placeins,needspace}
\usepackage{fvextra}
\fvset{fontsize=\footnotesize,baselinestretch=1,breaklines,breakanywhere,breaksymbolleft={},breaksymbolright={},breakanywheresymbolpre={}}
\definecolor{reportblue}{HTML}{173B57}
\usepackage{url}
\usepackage{hyperref}
\hypersetup{colorlinks=true,urlcolor=reportblue,linkcolor=black,unicode=true,hypertexnames=false,pdftitle={香山昆明湖V2自定义向量点积指令：测试数据与运行结果},pdfsubject={三阶段报告的测试数据与原始运行记录}}
\urlstyle{tt}
\usepackage{fancyhdr}
\pagestyle{fancy}\fancyhf{}
\fancyhead[L]{\small 香山昆明湖 V2 · 自定义向量点积指令}
\fancyhead[R]{\small 测试数据与运行结果}
\fancyfoot[C]{\thepage}
\ctexset{section={format=\Large\bfseries\rmfamily},subsection={format=\large\bfseries\rmfamily}}
\setlength{\parindent}{2em}
\setlength{\parskip}{3pt}
\linespread{1.22}
\setlength{\LTpre}{6pt}\setlength{\LTpost}{6pt}
\setlength{\LTcapwidth}{\textwidth}
\renewcommand{\arraystretch}{1.2}
\emergencystretch=3em
\setcounter{tocdepth}{2}
\captionsetup[figure]{font=small,labelsep=quad,justification=centering,position=bottom}
\captionsetup[table]{font=small,labelsep=quad,justification=centering,position=top}
\renewcommand{\topfraction}{0.9}\renewcommand{\bottomfraction}{0.6}
\renewcommand{\textfraction}{0.08}\renewcommand{\floatpagefraction}{0.75}
\graphicspath{{../figures/}{../evidence/G01_report-screenshots/}}
% 本材料内的文件路径，相对于材料根目录。
\newcommand{\ev}[1]{\texttt{\expandafter\seqsplit\expandafter{\detokenize{#1}}}}
% 团队仓库固定提交中的文件。
\newcommand{\repo}[3]{\href{https://github.com/cc77-builder/riscv/blob/#1/#2}{#3}}
\newcommand{\Rel}{f985cfecfa3b3760286d532d142878fd63cca6f1}
